\documentclass[11pt]{article}
\usepackage[margin=2.4cm]{geometry}
\usepackage{fontspec}
\setmainfont{DejaVu Sans}
\usepackage{juliamono}
\usepackage{fvextra}
\usepackage{xcolor}
\usepackage{xurl}
\usepackage{hyperref}
\usepackage{amssymb}
\newcommand{\checkbox}{$\square$}
\newcommand{\answerchoices}{\checkbox\ YES \quad \checkbox\ NO \quad \checkbox\ CANNOT DETERMINE}
\newcommand{\decisionexplanation}{%
  \par\smallskip\textbf{If NO, give the reason for the mismatch (required). If
  CANNOT DETERMINE, state what context or evidence is missing:}\par
  \rule{\linewidth}{0.4pt}\par\vspace{7pt}\rule{\linewidth}{0.4pt}\par}
\setlength{\parindent}{0pt}\setlength{\parskip}{5pt}
\setlength{\emergencystretch}{2em}
\DefineVerbatimEnvironment{Vsrc}{Verbatim}{fontsize=\small,frame=single,
  breaklines=true,breakanywhere=true}
\DefineVerbatimEnvironment{Vlean}{Verbatim}{fontsize=\scriptsize,frame=single,
  numbers=left,numbersep=5pt,breaklines=true,breakanywhere=true}
\hypersetup{colorlinks=true,urlcolor=blue}

\begin{document}
\begin{center}{\LARGE Proof Verification Document}\\[2pt]{\large reddit\_askmath\_pe1euu}\\[2pt]{\small Reddit proof-check posts (newly confirmed)}\end{center}
\begin{center}\fcolorbox{blue!60!black}{blue!4}{\parbox{0.92\linewidth}{\textbf{Make three independent decisions.} Follow the evidence chain: original source $\longrightarrow$ challenged claim $\longrightarrow$ proposed counterexample $\longrightarrow$ Lean theorem. A record is valid only when all three links are correct. Complete the human sign-off before reading the automated review appendix; those earlier results are background, not an answer key.}}\end{center}
\textbf{Use only this document, except for the public source link.} Everything needed for Questions 2 and 3 is included here. For Question 1, open the public source link when one is shown. If there is no public link, the complete stored problem and solution are included here. If the information is still not enough, choose CANNOT DETERMINE.
For each question, choose \textbf{YES}, \textbf{NO}, or \textbf{CANNOT DETERMINE}. Choose CANNOT DETERMINE if the document does not provide enough information for you to decide; do not mark YES if you are uncertain. If you find a specific mismatch, choose NO and explain it in the space below that question. Do not rewrite the claim, add assumptions, or substitute a different counterexample.
\textbf{How copied text is shown.} The problem, quoted passage, and proposed counterexample come from the stored record. The Lean theorem statement and Lean files are copied exactly from the certificate. Some long lines continue on the next printed line; this is only page formatting and does not change the copied text. Line numbers are added only to help you find passages. If the stored record contains a character that cannot be printed, the document shows \texttt{[U+XXXX CONTROL]} in its place.
\begin{description}
\item[Challenged claim.] The exact statement from the source that this record says is false.
\item[Proposed counterexample.] A concrete case that should satisfy the claim's assumptions but make its conclusion false.
\item[Lean certificate.] A Lean file proving a formal theorem. Successful compilation proves the theorem as written, not that it is the right translation of the source.
\end{description}
\section*{Question 1 --- Does this proof verification document accurately represent the source?}
\textbf{What to decide.} Compare the challenged passage below with the original source---either the public link or the complete source text in this document. Answer YES only if they make the same claim and include the same assumptions or conditions. If this document changes the meaning, leaves out a condition, or points to a different passage, answer NO.
\subsection*{Original source}
\textbf{Source record:}
\begin{Vsrc}
qedbench_certificate_availability_triage/data/records_reddit_exhaustive.jsonl
\end{Vsrc}
\textbf{Public source:}\par
\url{https://www.reddit.com/r/askmath/comments/pe1euu/could_someone_verify_my_graph_theory_proof/}\par
\subsection*{Problem or question from the source (verbatim)}
\begin{Vsrc}
Could someone verify my graph theory proof?
\end{Vsrc}
\subsection*{Challenged passage recorded from the source (stored text)}
\begin{Vsrc}
Let u be the uncolored vertex, if by coloring u we don't create a cycle, then G is a tree and therefore 2-colorable.
\end{Vsrc}
\subsection*{Surrounding source lines (stored text with added line labels)}
The labels refer to lines in the solution text stored with the source record; the source-line content is not transliterated.
\begin{Vsrc}
0001 | The question is "find if a graph G is 2-colorable".
0002 | 
0003 | I am trying to show that if a G has cycles, they need to be of even length in order to be 2-colorable.
0004 | 
0005 | Let be G a connected undirected graph with size |V| = k.
0006 | 
0007 | With induction, for k = 1, the statement holds.
0008 | 
0009 | Assume it holds for k = m, for k = m+1:
0010 | 
0011 | Let u be the uncolored vertex, if by coloring u we don't create a cycle, then G is a tree and therefore 2-colorable.
0012 | 
0013 | If by coloring u we create a cycle, then trivially if the cycle has an even length it is 2-colorable, otherwise it is not.
0014 | 
0015 | &#x200B;
0016 | 
0017 | I am not sure if the jump from "the cycle is 2-colorable" to "therefore the graph is 2-colorable" is valid.
0018 | 
0019 | Thanks in advance.
\end{Vsrc}
\textbf{A YES means:} the proof verification document faithfully represents both the source context and the challenged claim, without dropping a condition or changing its meaning.
\section*{Question 2 --- Does the proposed counterexample disprove that claim?}
\textbf{Claim to test (repeated for convenience).}
\begin{Vsrc}
Let u be the uncolored vertex, if by coloring u we don't create a cycle, then G is a tree and therefore 2-colorable.
\end{Vsrc}
\subsection*{How the record says the argument uses this claim}
This is the record's description, not an established fact; check it against the source.
\begin{Vsrc}
If adding/coloring the remaining vertex u creates no cycle, then the resulting connected graph G is a tree and therefore 2-colorable.
\end{Vsrc}
\subsection*{Proposed counterexample (exact values or construction)}
\begin{Vsrc}
Let V={a,b,c,u} and E={{a,b},{b,c},{c,a},{a,u}}. Before adding u, the graph is the triangle on {a,b,c}. Adding u with only edge {a,u} creates no new cycle containing u, since deg(u)=1. The resulting G is connected. Nevertheless, a-b-c-a is a cycle, so G is not a tree. It is not 2-colorable: a and b must differ, b and c must differ, hence a and c agree, contradicting edge {c,a}.
\end{Vsrc}
\subsection*{Why the record says the counterexample works}
Check this explanation independently; it is evidence to assess, not an answer key.
\begin{Vsrc}
The premise excludes a cycle created through u, but does not exclude a cycle already present among the other vertices. Thus the claimed tree and 2-colorability conclusions both fail.
\end{Vsrc}
\subsection*{How to decide}
\begin{enumerate}
\item List every assumption that applies to the challenged claim.
\item Check that the proposed example satisfies every one of those assumptions.
\item Recompute the claimed outcome and the actual outcome using exact reasoning.
\item Confirm that the conflict is with the challenged claim itself, not an unrelated sentence.
\end{enumerate}
\textbf{A YES means:} the example is fully specified, satisfies all relevant assumptions, and makes the challenged claim false. If a necessary value, condition, or calculation is missing, choose CANNOT DETERMINE.
\section*{Question 3 --- Does the final Lean theorem statement verify this same counterexample?}
\textbf{What to decide.} A Lean file can compile while proving the wrong theorem. Compare these three links independently:
\begin{enumerate}
\item The Lean definitions represent the same mathematical objects, assumptions, and operations as the challenged claim.
\item The Lean construction uses the same concrete values or example shown in Question 2.
\item The final theorem statement asserts that the challenged claim fails on that example, not merely a weaker or nearby statement.
\end{enumerate}
\subsection*{Selected final Lean theorem statement (verbatim UTF-8 excerpt)}
This is an exact contiguous slice of the certificate file, from the theorem keyword through the original \texttt{:=} token where its proof begins. It is shown first so the formal target is not buried in the complete file.
\begin{Vlean}
theorem reddit_askmath_pe1euu_refutation :
    simpleUndirected finalAdj = true /\
    simpleUndirected predecessorAdj = true /\
    predecessorHasExactlyABC = true /\
    predecessorIsInduced = true /\
    finalIsPredecessorWithLeafU = true /\
    connected finalAdj = true /\
    uHasExactlyNeighborA = true /\
    hasSimpleCycleThrough finalAdj .u = false /\
    simpleCycle predecessorAdj [.a, .b, .c, .a] = true /\
    hasSimpleCycle predecessorAdj = true /\
    acyclic predecessorAdj = false /\
    twoColorable predecessorAdj = false /\
    simpleCycle finalAdj [.a, .b, .c, .a] = true /\
    hasSimpleCycle finalAdj = true /\
    acyclic finalAdj = false /\
    tree finalAdj = false /\
    twoColorable finalAdj = false /\
    negativeControlSummary = true /\
    Not PublishedClaimAtWitness :=
\end{Vlean}
\subsection*{Project-local imported Lean files (complete, verbatim UTF-8)}
This section follows the certificate's import commands recursively and prints every repository-owned Lean source file on which the final theorem statement depends. Each file is copied in full, including its definitions and controls. External Mathlib, Std, and Lean Init source is not duplicated.
No project-local imported Lean files were found. External Mathlib, Std, and Lean Init source is not duplicated in this proof verification document.
\subsection*{Complete counterexample Lean file (verbatim UTF-8)}
The complete Lean file is included below so you can check every definition and supporting result used by the final theorem statement. Use it together with the project-local imported Lean files above.
\textbf{Lean certificate file:}
\begin{Vsrc}
partB_feasibility/reddit_certificates/reddit_askmath_pe1euu.lean
\end{Vsrc}
\begin{Vlean}
/-!
Record: reddit_askmath_pe1euu
Public source: https://www.reddit.com/r/askmath/comments/pe1euu/could_someone_verify_my_graph_theory_proof/

Exact target scope: the quoted induction branch claims that, when adjoining the
remaining vertex `u` creates no new cycle, the resulting connected graph is a
tree and hence is two-colourable.  The concrete witness has vertices
`{a,b,c,u}` and exactly the undirected edges `a-b`, `b-c`, `c-a`, and `a-u`.
The predecessor induced on exactly `{a,b,c}` is the triangle, and the final
graph is obtained by adjoining `u` with sole neighbour `a`.  Thus no simple
cycle contains `u`, while the old closed sequence `a,b,c,a` remains a simple
cycle; both the predecessor and final graph are not two-colourable.

Material modelling judgments: adjacency is an exact Boolean relation and is
separately checked symmetric and irreflexive.  Walks and paths are lists of
adjacent vertices, with paths additionally having no repeated vertex.  A simple
cycle is a closed adjacent list with at least three edges whose vertices before
the repeated endpoint are distinct.  Since this type has exactly four
vertices, exhaustive closed lists with three or four distinct vertices are
exactly all possible simple cycles; `acyclic` is their absence.  Connectivity
exhausts all simple paths of at most four vertices, which is exact on this
finite type.  Proper two-colourability exhausts every function from the four
vertices to `Bool`.  “Creates no cycle” is interpreted, as in the full-source
semantic audit, as absence of a new simple cycle through the adjoined `u`, not
as the already-false assertion that the whole predecessor is acyclic.

Negative control: deleting exactly `c-a` leaves precisely `a-b`, `b-c`, and
`a-u`.  Exhaustive cycle checking proves this connected graph acyclic/a tree,
and the explicit colouring `a=c=false`, `b=u=true` is proper.  This detects a
vacuous or uniformly-false encoding and detects omission of the deleted edge.
-/

namespace RedditAskmathPe1euu

inductive Vertex where
  | a | b | c | u
  deriving DecidableEq, BEq

abbrev Adj := Vertex -> Vertex -> Bool

def allVertices (p : Vertex -> Bool) : Bool :=
  p .a && p .b && p .c && p .u

def anyVertex (p : Vertex -> Bool) : Bool :=
  p .a || p .b || p .c || p .u

def finalAdj : Adj
  | .a, .b | .b, .a => true
  | .b, .c | .c, .b => true
  | .c, .a | .a, .c => true
  | .a, .u | .u, .a => true
  | _, _ => false

def inPredecessor : Vertex -> Bool
  | .a | .b | .c => true
  | .u => false

def predecessorAdj (x y : Vertex) : Bool :=
  inPredecessor x && inPredecessor y && finalAdj x y

def leafEdge : Adj
  | .a, .u | .u, .a => true
  | _, _ => false

def controlAdj : Adj
  | .a, .b | .b, .a => true
  | .b, .c | .c, .b => true
  | .a, .u | .u, .a => true
  | _, _ => false

def sameAdj (r s : Adj) : Bool :=
  allVertices fun x => allVertices fun y => (r x y) == (s x y)

def simpleUndirected (adj : Adj) : Bool :=
  (allVertices fun x => !(adj x x)) &&
  (allVertices fun x => allVertices fun y =>
    (adj x y) == (adj y x))

def containsVertex (needle : Vertex) : List Vertex -> Bool
  | [] => false
  | x :: xs => (needle == x) || containsVertex needle xs

def noRepeatedVertices : List Vertex -> Bool
  | [] => true
  | x :: xs => !(containsVertex x xs) && noRepeatedVertices xs

def walk (adj : Adj) : List Vertex -> Bool
  | [] => false
  | [_] => true
  | x :: y :: rest => adj x y && walk adj (y :: rest)

def path (adj : Adj) (vertices : List Vertex) : Bool :=
  walk adj vertices && noRepeatedVertices vertices

def initVertices : List Vertex -> List Vertex
  | [] => []
  | [_] => []
  | x :: xs => x :: initVertices xs

def atLeastThreeEdges : List Vertex -> Bool
  | _ :: _ :: _ :: _ :: _ => true
  | _ => false

def startsAndEndsSame : List Vertex -> Bool
  | [] => false
  | first :: rest =>
      let rec lastEquals : List Vertex -> Bool
        | [] => true
        | [last] => last == first
        | _ :: xs => lastEquals xs
      lastEquals rest

def simpleCycle (adj : Adj) (vertices : List Vertex) : Bool :=
  atLeastThreeEdges vertices &&
  startsAndEndsSame vertices &&
  walk adj vertices &&
  noRepeatedVertices (initVertices vertices)

def hasSimpleCycle (adj : Adj) : Bool :=
  (anyVertex fun x => anyVertex fun y => anyVertex fun z =>
    simpleCycle adj [x, y, z, x]) ||
  (anyVertex fun w => anyVertex fun x => anyVertex fun y => anyVertex fun z =>
    simpleCycle adj [w, x, y, z, w])

def hasSimpleCycleThrough (adj : Adj) (marked : Vertex) : Bool :=
  (anyVertex fun x => anyVertex fun y => anyVertex fun z =>
    simpleCycle adj [x, y, z, x] &&
    containsVertex marked (initVertices [x, y, z, x])) ||
  (anyVertex fun w => anyVertex fun x => anyVertex fun y => anyVertex fun z =>
    simpleCycle adj [w, x, y, z, w] &&
    containsVertex marked (initVertices [w, x, y, z, w]))

def acyclic (adj : Adj) : Bool :=
  !(hasSimpleCycle adj)

def hasPath (adj : Adj) (source target : Vertex) : Bool :=
  ((source == target) && path adj [source]) ||
  path adj [source, target] ||
  (anyVertex fun x => path adj [source, x, target]) ||
  (anyVertex fun x => anyVertex fun y =>
    path adj [source, x, y, target])

def connected (adj : Adj) : Bool :=
  allVertices fun source => allVertices fun target =>
    hasPath adj source target

def tree (adj : Adj) : Bool :=
  connected adj && acyclic adj

def coloring (ca cb cc cu : Bool) : Vertex -> Bool
  | .a => ca
  | .b => cb
  | .c => cc
  | .u => cu

def properColoring (adj : Adj) (color : Vertex -> Bool) : Bool :=
  allVertices fun x => allVertices fun y =>
    !(adj x y) || !(color x == color y)

def anyBool (p : Bool -> Bool) : Bool :=
  p false || p true

def twoColorable (adj : Adj) : Bool :=
  anyBool fun ca => anyBool fun cb => anyBool fun cc => anyBool fun cu =>
    properColoring adj (coloring ca cb cc cu)

def predecessorIsInduced : Bool :=
  allVertices fun x => allVertices fun y =>
    predecessorAdj x y ==
      (inPredecessor x && inPredecessor y && finalAdj x y)

def predecessorHasExactlyABC : Bool :=
  inPredecessor .a && inPredecessor .b && inPredecessor .c &&
  !(inPredecessor .u)

def finalIsPredecessorWithLeafU : Bool :=
  sameAdj finalAdj (fun x y => predecessorAdj x y || leafEdge x y)

def uHasExactlyNeighborA : Bool :=
  allVertices fun x => finalAdj .u x == (x == .a)

def negativeControlSummary : Bool :=
  simpleUndirected controlAdj &&
  connected controlAdj &&
  acyclic controlAdj &&
  tree controlAdj &&
  properColoring controlAdj (coloring false true false true) &&
  twoColorable controlAdj

abbrev PublishedClaimAtWitness : Prop :=
  simpleUndirected finalAdj = true ->
  simpleUndirected predecessorAdj = true ->
  predecessorHasExactlyABC = true ->
  predecessorIsInduced = true ->
  finalIsPredecessorWithLeafU = true ->
  connected finalAdj = true ->
  uHasExactlyNeighborA = true ->
  hasSimpleCycleThrough finalAdj .u = false ->
  simpleCycle predecessorAdj [.a, .b, .c, .a] = true ->
  twoColorable predecessorAdj = false ->
  tree finalAdj = true /\ twoColorable finalAdj = true

/-- The exact sole-neighbour premise, together with simplicity, rules out every
exhaustively represented simple cycle through the newly adjoined vertex `u`. -/
theorem no_new_simple_cycle_through_u_from_unique_neighbor :
    simpleUndirected finalAdj = true ->
    uHasExactlyNeighborA = true ->
    hasSimpleCycleThrough finalAdj .u = false := by
  decide

/-- The triangle defeats every proposed Boolean colouring, stated for an
arbitrary colour function rather than only as a computed existential flag. -/
theorem final_triangle_defeats_every_coloring (color : Vertex -> Bool) :
    properColoring finalAdj color = false := by
  cases ha : color .a <;>
  cases hb : color .b <;>
  cases hc : color .c <;>
  simp [properColoring, allVertices, finalAdj, ha, hb, hc]

/-- The induction-stage predecessor is itself the odd triangle, so its
non-two-colourability is preserved rather than silently replaced by a stronger
induction hypothesis. -/
theorem predecessor_triangle_defeats_every_coloring
    (color : Vertex -> Bool) :
    properColoring predecessorAdj color = false := by
  cases ha : color .a <;>
  cases hb : color .b <;>
  cases hc : color .c <;>
  simp [properColoring, allVertices, predecessorAdj, inPredecessor,
    finalAdj, ha, hb, hc]

/-- Load-bearing refutation certificate.  All source-context premises hold:
the predecessor is the induced odd triangle, the final graph is its connected
simple extension by leaf `u`, and no new simple cycle goes through `u`.  The
old triangle nevertheless makes the final graph cyclic, not a tree, and not
two-colourable.  The derived negative-control summary is true as well. -/
theorem reddit_askmath_pe1euu_refutation :
    simpleUndirected finalAdj = true /\
    simpleUndirected predecessorAdj = true /\
    predecessorHasExactlyABC = true /\
    predecessorIsInduced = true /\
    finalIsPredecessorWithLeafU = true /\
    connected finalAdj = true /\
    uHasExactlyNeighborA = true /\
    hasSimpleCycleThrough finalAdj .u = false /\
    simpleCycle predecessorAdj [.a, .b, .c, .a] = true /\
    hasSimpleCycle predecessorAdj = true /\
    acyclic predecessorAdj = false /\
    twoColorable predecessorAdj = false /\
    simpleCycle finalAdj [.a, .b, .c, .a] = true /\
    hasSimpleCycle finalAdj = true /\
    acyclic finalAdj = false /\
    tree finalAdj = false /\
    twoColorable finalAdj = false /\
    negativeControlSummary = true /\
    Not PublishedClaimAtWitness := by
  decide

/-- Meaningful edge-deletion negative control.  Exactly `c-a` is absent while
all three retained undirected edges are present; the result is connected,
acyclic, a tree, and accepts the independently verified proper colouring. -/
theorem reddit_askmath_pe1euu_negative_control :
    controlAdj .c .a = false /\
    controlAdj .a .c = false /\
    controlAdj .a .b = true /\ controlAdj .b .a = true /\
    controlAdj .b .c = true /\ controlAdj .c .b = true /\
    controlAdj .a .u = true /\ controlAdj .u .a = true /\
    simpleUndirected controlAdj = true /\
    connected controlAdj = true /\
    hasSimpleCycle controlAdj = false /\
    acyclic controlAdj = true /\
    tree controlAdj = true /\
    properColoring controlAdj (coloring false true false true) = true /\
    twoColorable controlAdj = true := by
  decide

end RedditAskmathPe1euu

\end{Vlean}
\subsection*{Proof-safety information}
The automated scan found no unfinished proofs, new axiom declarations, unsafe declarations, external code hooks, or uses of Lean's compiled decision procedure. This scan concerns proof construction only; it does not establish that the theorem matches the source claim.
\textit{Technical scan counts:} native\_decide=0, sorry=0, axiom\_decl=0, unsafe=0, extern=0. \textit{Technical tier:} KERNEL-ONLY.
\textbf{A YES means:} Lean represents the same objects and assumptions, uses the same proposed counterexample, and the final theorem statement asserts the same contradiction. Successful compilation by itself is not enough.
\clearpage
\section*{Human sign-off}
Complete this block from your own review of Questions 1--3. Do not use the automated appendix below as an answer key.\par
\textbf{Reviewer name:} \rule{9cm}{0.4pt}\par
\textbf{Date:} \rule{4cm}{0.4pt}\par
\medskip\textbf{1. Does the proof verification document accurately represent the source and challenged claim?}\par
\answerchoices\par
\decisionexplanation
\medskip\textbf{2. Does the proposed counterexample satisfy the assumptions and make that claim false?}\par
\answerchoices\par
\decisionexplanation
\medskip\textbf{3. Does the final Lean theorem statement represent and verify that same counterexample?}\par
\answerchoices\par
\decisionexplanation
\medskip\textbf{Overall decision:}\par
\checkbox\ VALIDATED (all three answers are YES)\par
\checkbox\ REJECTED (at least one answer is NO)\par
\checkbox\ NEEDS ANOTHER REVIEWER (at least one CANNOT DETERMINE and no NO)\par
\medskip\textbf{Additional comments.} Use this space for any extra source, counterexample, calculation, or Lean details that support your decisions.\par\rule{\linewidth}{0.4pt}\par\vspace{7pt}\rule{\linewidth}{0.4pt}
\clearpage
\section*{Appendix --- Prior automated review (background only)}
\textbf{These results are not an answer key.} They come from the earlier Opus/Fable verification campaign and are printed only so an auditor can see what was previously checked. The human decisions above must be made independently. Raw status codes are retained in small type for traceability.
\textbf{Overall automated disposition:} Cleared by automated screening; an independent search also found a counterexample.
\par\textit{Audit code:} \nolinkurl{PASS_CORROBORATED}
\subsection*{Independent search with no access to the stored example}
\textbf{Result:} An independent search found a counterexample without seeing the stored example.
\par\textit{Audit code:} \nolinkurl{WITNESS_FOUND}
\par\textbf{Cases checked:}
\begin{Vsrc}
Exhaustive over every labelled simple graph on n vertices, n = 1..7, with the uncoloured vertex fixed at label 0 (WLOG, since all labelled graphs on {0,...,n-1} are enumerated, so every isomorphism type of the pair (G,u) occurs). 2131019 graphs swept = 1+2+8+64+1024+32768+2097152. Each graph classified under all three antecedent readings plus a control.
\end{Vsrc}
\subsection*{Independent re-check of the proposed counterexample}
\textbf{Result:} A separate checker recomputed the proposed counterexample and confirmed it.
\par\textit{Audit code:} \nolinkurl{CONFIRMED}
\par\textbf{Checker's note:}
\begin{Vsrc}
Independent recomputation reproduces the refutation exactly. On G = triangle {a,b,c} plus pendant edge {a,u}, the vertex u has degree 1, lies on no simple cycle, and cycles(G) = cycles(G-u) = {abc}, so the quoted premise ('by coloring u we don't create a cycle') is TRUE under both the strong reading (no new cycle) and the local reading (u on no cycle), and G is connected. Yet BOTH halves of the conclusion are false: G has 4 vertices and 4 edges with the cycle a-b-c-a, so it is not a tree; (Shortened here; the complete note remains in the machine-check JSON.)
\end{Vsrc}
\subsection*{Independent review of the source-to-Lean correspondence}
\textbf{Result:} A separate reviewer judged that Lean proves the same failure using the same example.
\par\textit{Audit code:} \nolinkurl{MATCHES}
\par\textbf{Theorem summary from that reviewer:}
\begin{Vsrc}
The load-bearing theorem `reddit_askmath_pe1euu_refutation` is a single `decide`-proved conjunction over the fixed 4-vertex type {a,b,c,u} with `finalAdj` = exactly the undirected edges a-b, b-c, c-a, a-u. It asserts that every premise of the quoted induction branch holds at this instance -- `finalAdj` and the induced predecessor `predecessorAdj` are irreflexive and symmetric, the predecessor is exactly the induced subgraph on {a,b,c}, the final graph is the predecessor plus the single leaf. (Shortened here; the complete note remains in the machine-check JSON.)
\end{Vsrc}
\par\textbf{Difference noted by that reviewer:}
\begin{Vsrc}
None material; all three links hold. (1) Definitions: `simpleUndirected` checks irreflexivity and symmetry; walks are lists of pairwise-adjacent vertices; `simpleCycle` requires >=3 edges, closed, a walk, with all vertices before the repeated endpoint distinct -- the standard notion. `hasSimpleCycle` exhausts closed lists on 3 and 4 distinct vertices, which on a 4-element vertex type is exactly all simple cycles, so `acyclic` is exact rather than an under-approximation. (Shortened here; the complete note remains in the machine-check JSON.)
\end{Vsrc}
\subsection*{Follow-up review}
No follow-up review was needed or recorded.
\textit{This appendix is only background. You do not need any separate machine-output files to answer the three questions.}
\end{document}